\section{Sparse determinants with linearly many Boolean wires}
\label{sec:determinants}

The determinant will be a nonnegative connectivity weight after its
permutation terms have been summed. The required tensor representation
must therefore preserve permutation signs as well as matrix entries.
We give the sign construction explicitly; the inverse-pair accounting
is related to the determinant method in~\cite[Section 4]{BCKN15}.

\subsection{A sparse permutation-sign gadget}

Let $E\subseteq[N]\times[N]$ be the potential nonzero positions of
an $N\times N$ matrix, let $M=|E|$, and set
\[
 T_{\mathrm{det}}=N+M+2.
\]
Position bits $z_{ij}$ select one entry per row and column. A legal
selection is a permutation $\sigma$ supported on $E$. We construct
its sign with $O(T_{\mathrm{det}})$ edges and maximum degree
$O(T_{\mathrm{det}}^{2/3})$. The following construction is a
structural identity used under our standing assumption \HTC.

Put $B=\max\{1,\lfloor T_{\mathrm{det}}^{2/3}\rfloor\}$. Rows
having more than $B$ potential positions are heavy; each is a singleton
block. Split every consecutive run of light rows greedily, preserving
row order, so that the total number of potential positions in a block
is at most $B$. Columns are treated independently in their own order.
A zero potential row or column makes the determinant identically zero
and can be handled directly.

If there are $a$ row blocks and $b'$ column blocks, then
\begin{equation}\label{eq:block-count}
 a,b'=O(M/B+1),\qquad ab'=O(T_{\mathrm{det}}^{2/3}).
\end{equation}
Indeed there are at most $M/B$ heavy rows. In one light run, any two
consecutive greedy blocks have total load greater than $B$, so its
block count is at most twice the load divided by $B$, plus one.
There are at most one more light runs than heavy rows. In particular
$a\le4M/B+1$, and the column bound is identical. In a cell
$E_{uv}=E\cap(R_u\times C_v)$, a light side bounds its size by $B$;
if both sides are heavy singletons it contains at most one position.

For a legal permutation write $r(i)$ for its row-block index,
$c(j)$ for its column-block index, and
$n_{uv}=|\{i\in R_u:\sigma(i)\in C_v\}|$. Define
\begin{align*}
 I_R&=|\{i<i':r(i)=r(i'),\ c(\sigma(i))>c(\sigma(i'))\}|,\\
 I_C&=|\{j<j':c(j)=c(j'),\ r(\sigma^{-1}(j))>r(\sigma^{-1}(j'))\}|,\\
 I_{uv}&=|\{i<i':i,i'\in R_u,\ \sigma(i),\sigma(i')\in C_v,
                    \ \sigma(i)>\sigma(i')\}|,\\
 I_\times&=\sum_{u<u',\ v>v'}n_{uv}n_{u'v'}.
\end{align*}
The exact integer identity is
\begin{equation}\label{eq:sign-split}
 \inv(\sigma)=I_R+I_C+\sum_{u,v}I_{uv}+I_\times.
\end{equation}
To prove it, classify a pair $i<i'$ by whether its two row blocks
are equal and whether its two image column blocks are equal. If both
are equal, it is counted in $I_{uv}$. If only the row blocks are
equal, order preservation makes a column inversion equivalent to a
column-block inversion, counted in $I_R$. If only the column blocks
are equal, let $j<j'$ be the two image columns; an inversion is
exactly a reversal of the corresponding row-block order, counted in
$I_C$. If neither is equal, inversions occur precisely in cells with
$u<u'$ and $v>v'$, giving $I_\times$. The four classes are disjoint.

Use constant-state exact-one chains to check row and column selection.
Each light row block reads its at-most-$B$ position bits and contributes
$(-1)^{I_R}$ for that block; a heavy singleton contributes one.
Column blocks similarly contribute $(-1)^{I_C}$. Each cell reads its
position bits, contributes $(-1)^{I_{uv}}$, and outputs one parity bit
\[
 p_{uv}=\sum_{(i,j)\in E_{uv}}z_{ij}\pmod2.
\]
It can also reject repeated selected rows or columns in the cell.
Empty cells output a forced zero bit. A single factor reads the
$ab'$ parity bits and contributes
\begin{equation}\label{eq:cross-root}
 (-1)^{\sum_{u<u',\ v>v'}p_{uv}p_{u'v'}}.
\end{equation}
This gives $(-1)^{I_\times}$. Its degree is bounded by
\eqref{eq:block-count}. Every selected-position bit occurs in only a
constant number of checks. Total port count is
$O(N+M+ab')=O(T_{\mathrm{det}})$; maximum degree is
$O(B+ab')=O(T_{\mathrm{det}}^{2/3})$. Each entry is computed by
polynomially many arithmetic and pair-comparison operations on its
ports. Exact-one states and cell parities are unique on each legal
selection. Equality trees and subdivision preserve these bounds.

For each position multiply the factor
\[
 z_{ij}=0:\ 1,\qquad z_{ij}=1:\ A_{ij}.
\]
Summing over the position bits and the unique states gives precisely
\begin{equation}\label{eq:sparse-det}
 \sum_{\sigma\text{ supported on }E}
    \sgn(\sigma)\prod_iA_{i,\sigma(i)}=\det A.
\end{equation}
The $A_{ij}$ may depend on external Boolean variables. If those
variables and the total scopes of their conditions and numerical
entry factors are $O(T_{\mathrm{det}})$, the combined representation
still has linear edges. Its maximum degree is the larger of the
sign-gadget degree and the external degree. This is a representation
of a determinant inside a larger sum over external choices.

\subsection{Marker Laplacians and positivity}
\label{subsec:matrix-tree}

Fix a simple undirected graph $G=(V,E_G)$, $U\subseteq V$,
$F\subseteq E_G[U]$, and markers $M_0\subseteq U$. Define
\begin{equation}\label{eq:marker-matrix}
 A_{vv}=d_F(v)+[v\in M_0]+[v\notin U],\qquad
 A_{vw}=-[vw\in F]\quad(v\ne w).
\end{equation}
The unselected vertices form an identity block. On $U$, the matrix
is the Laplacian of the graph $H^+$ obtained from $H=(U,F)$ by
adjoining one root $\rho$ and the edges $\rho v$ for $v\in M_0$,
with the root row and column deleted.

For completeness, orient the edges of $H^+$ arbitrarily, and let
$B_\rho$ be its incidence matrix with the root row removed.
Then $A|_U=B_\rho B_\rho^T$, and Cauchy--Binet gives
\begin{equation}\label{eq:matrix-tree}
 \det A=\sum_{\substack{Q\subseteq E(H^+)\\|Q|=|U|}}
            \det((B_\rho)_Q)^2=\tau(H^+).
\end{equation}
A cyclic edge set has dependent incidence columns. An acyclic set
of $|U|$ edges on $|U|+1$ vertices is a spanning tree. Its incidence
minor has determinant $\pm1$, by repeatedly expanding at a nonroot
leaf. Thus every tree contributes one and every other set contributes
zero. Empty $U$ gives determinant one. In particular,
\begin{equation}\label{eq:marker-positive}
 \det A\ge0,\qquad \det A>0
 \ \Longleftrightarrow\
 \text{every component of }(U,F)\text{ meets }M_0.
\end{equation}
If $M_0=\{r\}$ with $r\in U$, positivity is connectivity of $(U,F)$,
and $\det A=\tau(U,F)$ because the root's single incident edge is a
bridge.

The potential positions in~\eqref{eq:marker-matrix} number at most
$n+2m$. For local evaluation, compute every $d_F(v)$ by an
$O(\deg_G(v)+1)$-gate Boolean population counter. One construction
repeatedly uses constant-size full adders to reduce triples in each
binary column: each full adder reduces the total number of live bits
by one, so $O(d)$ suffice for $d$ input bits, followed by addition
of two $O(\log(d+2))$-bit words. Its internal states are functional.
All degree circuits together have $O(n+m)$ wires and gates.
Their output-bit count is $O(n+m)$ as well. A diagonal numerical
factor reads only its degree output, marker bit, selected-vertex bit,
and position bit; an off-diagonal factor reads a position bit and an
edge-selection bit. Total scope length stays $O(n+m)$, including at
vertices of high graph degree.

\subsection{Exact statistics with long binary coefficients}
\label{subsec:weights}

The applications also use a fixed number of statistics
$A_j(y)=\sum_i a_{ji}y_i$ on $p=O(S)$ object bits, with nonnegative
coefficients of at most $H=\lceil S^{9/10}\rceil$ bits. Any fixed
polynomial-time relation on these statistics, including a threshold,
can be added using $O(S)$ wires, degree $O(S^{9/10})$, and a unique
auxiliary extension.

Indeed, group the object bits in blocks of
$B=\max\{1,\lfloor S^{9/10}\rfloor\}$ bits. There are
$h=O(S^{1/10}+1)$ blocks. Every partial or total sum fits in
\[
 q=H+\lceil\log_2(p+1)\rceil+1=O(S^{9/10})
\]
bits. A block factor reads its original bits and checks its output
partial sums. For each statistic, a binary tree of addition factors
combines these outputs. Each addition factor reads three $q$-bit
words, rejects overflow, and enforces exact addition; on genuine
inputs there is no overflow. The root relation reads only the fixed
number of total sums. The total number of state wires is $O(hq)=O(S)$,
whereas the maximum degree is $O(S^{9/10})$. The use of an addition
tree, rather than one factor reading every block sum, is essential to
this degree bound. Local entry arithmetic is polynomial in $S+L$.

For coefficients bounded in value by $S^{O(1)}$, a smaller-block
variant also works: take blocks of $S^{2/3}$ inputs, transmit
$O(\log S)$-bit partial sums, and let a terminal factor read all of
them. Its degree is $O(S^{1/3}\log S)$, and the block degree is
$O(S^{2/3})$. Both encodings are compatible with the determinant
representation. The larger-degree one uses the constant rescaling
\eqref{eq:rescale} before invoking \HTC.
